% 两个独立公平比特X,Y及其模二和Z=X xor Y。Omega={00,01,10,11}，每点概率1/4，格子(x,y)以第一位为列、第二位为行。
% (a)(b)(c)分别为sigma(X)、sigma(Y)、sigma(Z)的两个原子（各占两格，概率1/2）；
% 任两种划分相交得四个各1/4的格子，故两两独立；(d) sigma(X,Y)的四个原子中每格Z已确定，
% {X=1,Y=1,Z=1}为空集，概率0不等于1/8。
\begin{tikzpicture}[line width=.6pt,font=\fontsize{8}{10}\selectfont,>=Stealth]
\def\c{8mm}
\newcommand{\cells}[4]{% fills for 00,01,10,11 : (x,y) -> column x, row y
  \fill[#1] (0,0) rectangle (\c,\c);
  \fill[#2] (0,\c) rectangle (\c,2*\c);
  \fill[#3] (\c,0) rectangle (2*\c,\c);
  \fill[#4] (\c,\c) rectangle (2*\c,2*\c);
  \draw[FigureGrid,line width=.4pt] (\c,0) -- (\c,2*\c) (0,\c) -- (2*\c,\c);
  \draw[FigureStroke] (0,0) rectangle (2*\c,2*\c);
  \foreach \x/\y in {0/0,0/1,1/0,1/1}{\node[FigureInk] at ({(\x+.5)*\c},{(\y+.5)*\c}) {$\x\y$};}}
\begin{scope}
 \cells{FigureBlueFill}{FigureBlueFill}{FigureYellowFill}{FigureYellowFill}
 \draw[FigureStroke,line width=1.1pt] (\c,0) -- (\c,2*\c);
 \node[FigureInk] at (\c,-4mm) {(a) $\sigma(X)$};
\end{scope}
\begin{scope}[xshift=28mm]
 \cells{FigureBlueFill}{FigureYellowFill}{FigureBlueFill}{FigureYellowFill}
 \draw[FigureStroke,line width=1.1pt] (0,\c) -- (2*\c,\c);
 \node[FigureInk] at (\c,-4mm) {(b) $\sigma(Y)$};
\end{scope}
\begin{scope}[xshift=56mm]
 \cells{FigureBlueFill}{FigureYellowFill}{FigureYellowFill}{FigureBlueFill}
 \node[FigureInk] at (\c,-4mm) {(c) $\sigma(Z)$};
\end{scope}
\begin{scope}[xshift=84mm]
 \cells{FigureBlueFill}{FigureGreenFill}{FigureYellowFill}{FigurePurpleFill}
 \draw[FigureStroke,line width=1.1pt] (\c,0) -- (\c,2*\c) (0,\c) -- (2*\c,\c);
 \node[FigureInk] at (\c,-4mm) {(d) $\sigma(X,Y)$};
 \node[FigureMuted,font=\fontsize{7.5}{9}\selectfont] at (\c,-9mm) {每格内$Z$已确定};
\end{scope}
\node[FigureMuted,font=\fontsize{7.5}{9}\selectfont] at (64mm,-9mm) {蓝：$Z=0$；黄：$Z=1$};
\end{tikzpicture}
